\section{Proposed Solution}
\label{sec:solution}

We propose to recover skill lineage from the content of skills alone and then use it to
study reuse. The solution has two parts.

First, we build a classifier that decides, for a pair of skills, whether one is a derived
variant of the other, a template sibling, or unrelated. Separating these classes is the core
difficulty. Skills are short, mostly written in English, and often scaffolded from the same
templates, so two independently written skills can share much of their text while a reworded
or translated copy can share little. No fixed similarity threshold separates these cases, and
training a model on labels produced by such a threshold would only teach the threshold back to
the model. We therefore build the ground truth by hand: team members independently label real
candidate pairs from GitHub against a written rubric, and we measure their agreement. Because
the pairs are sampled across similarity levels, the labelled set deliberately contains the
highly similar template siblings that existing benchmarks leave out. Besides lexical
similarity, the classifier uses signals aimed at the template problem, chiefly a similarity
computed only over the lines a pair shares that do not recur across many unrelated skills,
together with front-matter, structural, and repository evidence. We compare it with exact
hashing, MinHash-style and TF-IDF similarity, and a re-implementation of
SkillReuse~\cite{zhu2026skillreuse}.

Second, we apply the classifier to candidate pairs drawn from the GitSkills
dataset~\cite{destefanis2027gitskills}, group skills into families, order family members by
the first commit of each \texttt{SKILL.md} where that history is available, and measure how
concentrated reuse is once edited copies are counted. In the next phase we will use these
families to test whether upstream changes reach their copies, following the genealogy view of
code clones~\cite{kim2005genealogy}.

The design answers each limitation identified above. It needs no registry entry and no fork
record, which matters in practice: GitHub code search indexes a fork only when it has more
stars than its parent, so only three repositories in GitSkills are
forks~\cite{gitskills_sample}. It separates derivation from template similarity instead of
treating every similar pair as a clone, and it tests that separation on real pairs rather than
synthetic mutations. It uses commit history to give derived pairs a direction where possible.
Finally, it forms families with community detection rather than plain connected components,
so that a few bridging pairs cannot chain unrelated skills into one giant family.
